\documentclass[11pt]{article}
\usepackage{amsmath,amssymb,amsthm}
\usepackage[margin=1in]{geometry}
\usepackage{hyperref}
\usepackage{enumitem}

\newtheorem{theorem}{Theorem}[section]
\newtheorem{lemma}[theorem]{Lemma}
\theoremstyle{definition}
\newtheorem{definition}[theorem]{Definition}
\theoremstyle{remark}
\newtheorem{remark}[theorem]{Remark}

\newcommand{\tagsource}{\textsf{\footnotesize[SOURCE]}}
\newcommand{\taglocal}{\textsf{\footnotesize[LOCAL]}}
\newcommand{\tagsynth}{\textsf{\footnotesize[SYNTHESIS]}}
\newcommand{\tagconj}{\textsf{\footnotesize[CONJECTURE]}}

\title{Finite $p$-Variation Jordan Curves and\\ Inscribed Rectangles via Asano--Ike}
\author{Internal proof package --- clean-room verified (Directive 007R)\\Not externally peer-reviewed}
\date{June 18, 2026}

\begin{document}
\maketitle

\begin{abstract}
We record an extension of the rectangular peg theorem to Jordan curves
with finite $p$-variation for $1\le p<2$, combining Boedihardjo--Geng
piecewise-linear interpolation \tagsource, a velocity-matched cubic
polygon-smoothing lemma T018 \taglocal, Young primitive convergence
\taglocal, and Asano--Ike Theorem~1.1 \tagsource.
Clean-room audit (Directive~007R) verified the proof chain.
This document is \textbf{internally} verified only.
We do \textbf{not} treat the critical $p=2$ boundary (see \S\ref{sec:limit}).
\end{abstract}

\section{Introduction}

The Toeplitz square peg problem asks whether every Jordan curve in $\mathbb{R}^2$
inscribes a square. Asano and Ike (arXiv:2412.21057v3) prove that if a Jordan curve
satisfies a continuous Legendrian-lift hypothesis (Theorem~1.1), then it inscribes
every $\theta$-rectangle for $\theta\in(0,\pi)$.

\begin{theorem}[T017 \tagsynth]\label{thm:T017}
Let $c:S^1\to\mathbb{R}^2$ be a Jordan parametrization with finite $p$-variation
for some $1\le p<2$. Then $c$ satisfies the hypotheses of Asano--Ike Theorem~1.1
and therefore inscribes a $\theta$-rectangle for every $\theta\in(0,\pi)$.
\end{theorem}

\begin{remark}[Novelty language]
This appears to be an immediate synthesis of existing interpolation and microlocal
peg results, with the controlled $C^1$ smoothing lemma supplied here.
We do not claim priority without author confirmation.
\end{remark}

\section{Variation preliminaries \taglocal}

\begin{definition}
For $u:[a,b]\to\mathbb{R}^d$,
\[
|u|_{r\text{-var};[a,b]}
=\Bigl(\sup_{\Pi}\sum_i|u(t_{i+1})-u(t_i)|^r\Bigr)^{1/r},
\qquad
\|u\|_\infty=\sup_t|u(t)|.
\]
Constant offsets are not folded into the seminorm.
\end{definition}

\begin{lemma}[Interpolation \taglocal]\label{lem:interp}
For $0<p<q$ and finite $p$-variation $h$,
\[
|h|_{q\text{-var}}\le \operatorname{osc}(h)^{1-p/q}|h|_{p\text{-var}}^{p/q},
\quad \operatorname{osc}(h)=\sup_{s,t}|u(t)-u(s)|.
\]
\end{lemma}

\begin{lemma}[Variation monotonicity \taglocal]\label{lem:qle1}
For $q\ge 1$, $|h|_{q\text{-var}}\le |h|_{1\text{-var}}$.
\end{lemma}

\begin{proof}
For each partition, $(\sum|\Delta h_i|^q)^{1/q}\le\sum|\Delta h_i|$ since
$x^q\le x$ for $x\in[0,1]$ applied after normalizing by $\max_i|\Delta h_i|$
when $\max_i|\Delta h_i|\le 1$; the general case uses
$(\sum x_i^q)^{1/q}\le\sum x_i$ for nonnegative $x_i$.
\end{proof}

\section{Polygonal interpolation \tagsource}

\begin{theorem}[Boedihardjo--Geng 2.2, $M=\mathbb{R}^2$ \tagsource]
For a Jordan curve $\gamma$ and $\varepsilon>0$, there is a partition
$\mathcal{P}_\varepsilon$ with mesh $<\varepsilon$ such that the piecewise
linear interpolant $\gamma^{\mathcal{P}_\varepsilon}$ is a simple closed polygon
with vertices on $\gamma$ in cyclic order on the same parameter domain.
\end{theorem}

\begin{lemma}[Boedihardjo--Geng 3.1, $f=\mathrm{id}$ \tagsource]
If $\gamma$ has finite $p$-variation and $q>p$, then
$\|\gamma-\gamma^{\mathcal{P}}\|_{q\text{-var}}\to 0$ as $|\mathcal{P}|\to 0$.
\end{lemma}

\section{T018: velocity-matched cubic smoothing \taglocal}

\begin{lemma}[T018 \taglocal]\label{lem:T018}
Let $P:S^1\to\mathbb{R}^2$ be a simple polygonal Jordan curve.
For every $\varepsilon>0$ and $q\ge 1$, there exists a regular $C^1$ Jordan
map $S$ with the same orientation and a common parametrization such that
$\|S-P\|_\infty<\varepsilon$ and $\|S-P\|_{q\text{-var}}<\varepsilon$.
\end{lemma}

\paragraph{Construction.}
At vertex $v=P(\tau)$, choose $\delta>0$, $A=P(\tau-\delta)$, $C=P(\tau+\delta)$,
$L=2\delta$, $u_-=(v-A)/\delta$, $u_+=(C-v)/\delta$.
Cubic controls:
\[
Q_0=A,\ Q_1=A+\tfrac23(v-A),\ Q_2=C-\tfrac23(C-v),\ Q_3=C.
\]
On $[\tau-\delta,\tau+\delta]$, $S$ is the cubic B\'ezier; elsewhere $S=P$.

\paragraph{Velocity matching.}
$\frac{dS}{dt}(\tau-\delta)=u_-$ and $\frac{dS}{dt}(\tau+\delta)=u_+$.

\paragraph{Injectivity.}
With $a=v-A$, $b=C-v$ non-antiparallel, $w=a/|a|+b/|b|$
satisfies $w\cdot dS/ds>0$ on the corner arc.

\paragraph{Variation.}
$\mathrm{Var}_1(S-P;I_v)\le 4r_v$ with $|a|,|b|\le r_v$; sum over corners.

\section{Young primitive convergence \taglocal}

Write $c=(x,\xi)$, $\lambda=\xi\,dx$, $f(t)=\int_0^t\xi\,dx$ (Young, $q<2$).

\begin{theorem}[Young / BG 3.1 \tagsource]
If $1/p+1/q>1$, the Young integral exists; for $q<2$, $2/q>1$.
\end{theorem}

With $\|S_n-c\|_{q\text{-var}}\to 0$,
\[
\sup_{t\in[0,2\pi]}|f_n(t)-f(t)|\le C_q\|S_n-c\|_{q\text{-var}}\to 0.
\]

\paragraph{Lift to $\mathbb{R}$ \taglocal.}
For $t=2\pi k+s$, $|k|\le N_K:=\lceil K/(2\pi)\rceil+1$,
\[
\sup_{|t|\le K}|f_n(t)-f(t)|\le \eta_n+N_K|I_n-I|\to 0.
\]

\section{Application of Asano--Ike \tagsource}

Normalize enclosed area to $\pi$. Theorem~1.1 (AI-PEG-v3) and Remark~5.6
($C^1$ approximants suffice) yield all $\theta$-rectangles.
Depends on preprint arXiv:2412.21057v3.

\section{Examples and limitations}\label{sec:limit}

The Koch snowflake is nonrectifiable with finite $p$-variation for
$p>\log 4/\log 3$. We do not claim it lies outside the locally monotone
class without a separate proof.

\paragraph{Unresolved $p=2$ boundary \tagconj.}
At $q=2$, Young requires $1/q+1/q=1$ (not strict).
Critical-modulus enhancement (Yang, arXiv:1204.0112) is treated separately
in Directive~008; T020 remains conjectural.

\bibliographystyle{plain}
\begin{thebibliography}{9}
\bibitem{AI2026}
T.~Asano and Y.~Ike,
\textit{The rectifiable rectangular peg problem},
arXiv:2412.21057v3, 2026.
\bibitem{BG2014}
H.~Boedihardjo and X.~Geng,
arXiv:1309.1576v2.
\bibitem{Yang2012}
D.~Yang,
\textit{Notes on area operator, geometric 2-rough paths and Young integral},
arXiv:1204.0112.
\end{thebibliography}

\end{document}
